\chapter{EML Registry: Tamari Lattice Contraction}

\section{Core Types}

\begin{definition}[EML Tree]
\label{def:emltree}
\lean{EMLRegistry.EMLTree}
\leanok
The type of binary trees representing configurations in the Tamari lattice.
A tree is either a \texttt{Leaf} (empty configuration) or a \texttt{Node} $t_1\ t_2$ (composition of two sub-configurations).
\end{definition}

\begin{definition}[Tree Size]
\label{def:size}
\lean{EMLRegistry.EMLTree.size}
\leanok
\uses{def:emltree}
The number of internal nodes in a tree, used as the lattice index $n$ in $T_n$.
\end{definition}

\begin{definition}[Single-Step Contraction]
\label{def:contracts_one}
\lean{EMLRegistry.contracts_one}
\leanok
\uses{def:emltree}
The primitive non-associative rewrite rule: $(a \bullet b) \bullet c \to a \bullet (b \bullet c)$.
This is the coupling signature for choice resolution, lifted through subtrees.
\end{definition}

\begin{definition}[Multi-Step Contraction]
\label{def:contracts_to}
\lean{EMLRegistry.contracts_to}
\leanok
\uses{def:contracts_one}
The reflexive-transitive closure of \textsf{contracts\_one}. Represents an audit trail of choice resolutions with monotonic provenance.
\end{definition}

\begin{definition}[Right-Comb Normal Form]
\label{def:rightComb}
\lean{EMLRegistry.rightComb}
\leanok
\uses{def:emltree}
The minimum element of the Tamari lattice — the equilibrium attractor.
Defined as $\mathsf{rightComb}(0) = \mathsf{Leaf}$ and $\mathsf{rightComb}(n+1) = \mathsf{Node}(\mathsf{Leaf},\ \mathsf{rightComb}(n))$.
\end{definition}

\section{Lifting Lemmas}

\begin{lemma}[Contraction Preserved Under Left Node]
\label{lem:node_left}
\lean{EMLRegistry.contracts_to_node_left}
\leanok
\uses{def:contracts_one, def:contracts_to}
If $t \to t'$ then $\mathsf{Node}(t,\ r) \to \mathsf{Node}(t',\ r)$.
Choice resolution in the left subtree propagates to the whole composition.
\end{lemma}

\begin{lemma}[Contraction Preserved Under Right Node]
\label{lem:node_right}
\lean{EMLRegistry.contracts_to_node_right}
\leanok
\uses{def:contracts_one, def:contracts_to}
If $r \to r'$ then $\mathsf{Node}(l,\ r) \to \mathsf{Node}(l,\ r')$.
Choice resolution in the right subtree propagates to the whole composition.
\end{lemma}

\begin{lemma}[Transitivity of Contraction]
\label{lem:trans}
\lean{EMLRegistry.contracts_to_trans}
\leanok
\uses{def:contracts_to}
If $s \to t$ and $t \to u$ then $s \to u$.
Audit trails concatenate monotonically.
\end{lemma}

\section{Size Invariants}

\begin{lemma}[Size is Preserved by Single-Step Contraction]
\label{lem:size_one}
\lean{EMLRegistry.contracts_one_size_eq}
\leanok
\uses{def:size, def:contracts_one}
Rotation rewrites the tree depth but does not change the number of internal nodes.
\end{lemma}

\begin{lemma}[Size is Preserved by Multi-Step Contraction]
\label{lem:size_multi}
\lean{EMLRegistry.contracts_to_size_eq}
\leanok
\uses{def:size, def:contracts_to, lem:size_one}
Contraction paths preserve tree size, establishing that the Tamari lattice $T_n$ is graded by node count $n$.
\end{lemma}

\section{Main Contraction Theorem}

\begin{lemma}[Node of Right-Combs Contracts to Right-Comb]
\label{lem:node_rightcombs}
\lean{EMLRegistry.node_of_rightCombs_contracts_to_rightComb}
\leanok
\uses{def:rightComb, def:contracts_one, def:contracts_to, lem:node_right}
For any $a, b : \mathbb{N}$,
\[
\mathsf{Node}(\mathsf{rightComb}(a),\ \mathsf{rightComb}(b)) \longrightarrow \mathsf{rightComb}(1 + a + b).
\]
The equilibrium attractor is closed under non-associative composition.
\end{lemma}

\begin{theorem}[Every Tree Contracts to Right-Comb]
\label{thm:contracts_to_rightComb}
\lean{EMLRegistry.contracts_to_rightComb}
\leanok
\uses{def:emltree, def:size, def:rightComb, def:contracts_to, lem:node_left, lem:node_right, lem:trans, lem:node_rightcombs}
For any tree $t : \mathsf{EMLTree}$,
\[
t \longrightarrow \mathsf{rightComb}(t.\mathsf{size}).
\]
Every configuration has a valid evolution path to the equilibrium attractor indexed by its size.
\end{theorem}

\section{Router and Registry}

\begin{definition}[Router Index]
\label{def:router}
\lean{EMLRegistry.RouterIndex}
\leanok
A bounded natural number $\mathsf{Fin}(n)$ serving as a neural binding address.
\end{definition}

\begin{definition}[Type Registry]
\label{def:registry}
\lean{EMLRegistry.TypeRegistry}
\leanok
\uses{def:router, def:emltree}
An injective mapping from router indices to EML trees. This is the cortex-registry interface where neural computation meets formal grammar.
\end{definition}

\section{Witness-Skeptic Game}

\begin{definition}[Cortex Certificate]
\label{def:certificate}
\lean{EMLRegistry.CortexCertificate}
\leanok
\uses{def:contracts_to, def:rightComb}
A proof-carrying audit trail $\langle source,\ target,\ proof \rangle$ showing that a tree reaches its equilibrium.
\end{definition}

\begin{definition}[Certify]
\label{def:certify}
\lean{EMLRegistry.certify}
\leanok
\uses{def:contracts_to, def:rightComb, def:certificate, thm:contracts_to_rightComb}
Issues a \textsf{CortexCertificate} for any tree $t$, certifying that $t \to \mathsf{rightComb}(t.\mathsf{size})$.
This is the quench witness in the Witness-Skeptic game.
\end{definition}

\chapter{Logic Types: Pluralistic Reasoning Framework}

\section{Logic Type Hierarchy}

\begin{definition}[13 Logic Types]
\label{def:logictype}
\lean{LogicTypes.LogicType}
\leanok
An inductive type enumerating 13 reasoning paradigms.

Each handles specific classes of paradoxes as boundary conditions rather than errors.
\end{definition}

\begin{definition}[Logic Classification]
\label{def:logicclass}
\lean{LogicTypes.LogicClass}
\leanok
\uses{def:logictype}
Classification of logic types as Substructural, Extension, Alternative, Generalization, or Classical.
\end{definition}

\begin{definition}[Logic Tree]
\label{def:logictree}
\lean{LogicTypes.LogicTree}
\leanok
\uses{def:emltree, def:logictype}
A type family assigning EML trees to each logic type. Currently uniform; extensible to logic-specific tree types.
\end{definition}

\begin{definition}[Logic Contraction]
\label{def:logiccontraction}
\lean{LogicTypes.LogicContraction}
\leanok
\uses{def:contracts_to, def:logictype}
A type family of contraction relations indexed by \textsf{LogicType}. For Classical logic this is the Tamari contraction from EMLRegistry.
\end{definition}

\begin{definition}[Logic Translation]
\label{def:translation}
\lean{LogicTypes.LogicTranslation}
\leanok
\uses{def:logictype, def:logiccontraction}
A structure specifying forward and backward maps between logic types with soundness and completeness proofs.
\end{definition}

\begin{definition}[Meta-Contraction]
\label{def:metacontracts}
\lean{LogicTypes.MetaContractsTo}
\leanok
\uses{def:logiccontraction, def:translation}
Cross-logic contraction allowing reasoning that spans multiple logic types via intra-logic contraction, inter-logic translation, and transitivity.
\end{definition}

\section{Cayley-Dickson Connection}

\begin{definition}[CD Step]
\label{def:cdstep}
\lean{LogicTypes.LogicType.cdStep}
\leanok
\uses{def:logictype}
Maps logic types to Cayley-Dickson construction steps.
\end{definition}

\begin{definition}[Split-Octonion Sectors]
\label{def:split}
\lean{LogicTypes.LogicType.isAssociativeSector}
\leanok
\uses{def:logictype}
Divides logic types into associative and non-associative sectors of the split-octonion algebra.
\end{definition}

\section{Normal Forms and Contraction Theorem}

\begin{definition}[Logic Normal Form]
\label{def:logicnormal}
\lean{LogicTypes.LogicNormalForm}
\leanok
\uses{def:rightComb, def:logictype}
The normal form for each logic type. For Classical logic this is \textsf{rightComb}.
\end{definition}

\begin{theorem}[Logic Contraction to Normal Form]
\label{thm:logic_contracts}
\lean{LogicTypes.logic_contracts_to_normal_form}
\leanok
\uses{def:logicnormal, def:logiccontraction, def:size, thm:contracts_to_rightComb}
For any logic type $lt$ and tree $t$, $\mathsf{LogicContraction}\ lt\ t\ (\mathsf{LogicNormalForm}\ lt\ t.\mathsf{size})$.
Proven for all 13 logic types via $\mathsf{contracts\_to\_rightComb}$.
\end{theorem}

\section{Logic Monad: Self-Contained Self-Reference}

\begin{definition}[Free Logic Monad]
\label{def:logic_monad}
\lean{LogicMonad.LogicM}
\leanok
\uses{def:emltree}
The free monad over binary trees: $\mathsf{LogicM}\ \alpha$ is a tree whose
leaves carry values of type $\alpha$ and whose internal nodes represent
non-associative composition of computations.
\begin{itemize}
  \item $\mathsf{pure}\ a$ = a leaf containing $a$ (a completed computation)
  \item $\mathsf{node}\ l\ r$ = deferred composition of two sub-computations
  \item $\mathsf{bind}$ = substitution: replace every leaf with a new tree
\end{itemize}
This is self-contained by definition (recursive structure IS self-reference)
and extensible (interpreters fold over $\mathsf{LogicM}\ \alpha$ to produce
any target monad).
\end{definition}

\begin{theorem}[Monad Laws for LogicM]
\label{thm:logic_monad_laws}
\lean{LogicMonad.pure_bind}
\lean{LogicMonad.bind_pure}
\lean{LogicMonad.bind_assoc}
\leanok
\uses{def:logic_monad}
The free monad $\mathsf{LogicM}$ satisfies the three monad laws:
\begin{itemize}
  \item $\mathsf{pure}\ a \gg\!\!= f = f a$ (left identity)
  \item $m \gg\!\!= \mathsf{pure} = m$ (right identity)
  \item $(m \gg\!\!= f) \gg\!\!= g = m \gg\!\!= (\lambda x.\ f x \gg\!\!= g)$ (associativity)
\end{itemize}
Proof by structural induction on the tree.
\end{theorem}

\begin{definition}[Logic-Specific Monad]
\label{def:logic_specific_monad}
\lean{LogicMonad.LogicMonad}
\leanok
\uses{def:logic_monad, def:logiccontraction, def:logicnormal}
For each logic type $\mathsf{lt}$, the $\mathsf{LogicMonad}\ \mathsf{lt}\ \alpha$
structure wraps a $\mathsf{LogicM}\ \alpha$ and normalizes it via
$\mathsf{LogicContraction}\ \mathsf{lt}$. The normalization cost is
$\mathsf{cdStep}\ \mathsf{lt}$.

This is the minimal monad for each logic: it does the least work
(pure substitution) and then normalizes using that logic's contraction
dynamics. The self-reference property of the free monad (bind = substitution)
means every logic monad shares the same core recursion — they differ only
in how they normalize.

The pluralistic logic monad \textbf{transcends identity}: the monad structure
(pure, bind, the laws) is invariant under change of logic. Only the
normalization differs. This means the ``eternal'' layer (pure functional
composition) is common to all logics; the ``endless'' layer (iteration,
normalization cost) is what distinguishes them.
\end{definition}

\begin{theorem}[Monad Structure is Invariant]
\label{thm:monad_invariant}
\lean{LogicMonad.monad_structure_invariant}
\leanok
\uses{def:logic_specific_monad, def:logic_monad}
The forgetful functor from $\mathsf{LogicMonad}\ \mathsf{lt}$ to $\mathsf{LogicM}$
is an isomorphism of monads for any $\mathsf{lt}$. Only the normalization
(\textsf{LogicContraction}) changes between logics.
\end{theorem}

\section{Institutional Closure: Interacting Logic Monads}

\begin{definition}[Institutional Closure]
\label{def:institutional_closure}
\lean{InstitutionalClosure.closure}
\leanok
\uses{def:logic_monad, def:emltree}
The institutional closure pipeline composes four layers:
\begin{enumerate}
  \item \textbf{Temporal normalization} --- linearize the event timeline via
    $\mathsf{LogicContraction\ Temporal\ =\ contracts\_to}$.
  \item \textbf{Fuzzy grading} --- evaluate each event's impact through
    graded membership (Sorites-like gradation).
  \item \textbf{Deontic update} --- revise norms from accumulated
    blame, using $\mathsf{LogicContraction\ Deontic}$.
  \item \textbf{Self-recognition} --- the monadic bind that identifies
    the output norm as the institution's own (fixed point).
\end{enumerate}
The pipeline is a monadic computation $\mathsf{LogicM\ Event \to LogicM\ Norm}$
using the shared $\mathsf{LogicM}$ backbone, with logic-specific normalization
at each layer.
\end{definition}

\begin{theorem}[Closure Fixed Point]
\label{thm:closure_fixed_point}
\lean{InstitutionalClosure.closure_is_fixed_point}
\leanok
\uses{def:institutional_closure}
$\mathsf{selfRecognize} \circ \mathsf{closure} = \mathsf{closure}$:
self-recognition is a fixed point. Once the pipeline has produced a norm,
further recognition does not change it.
\end{theorem}

\begin{theorem}[Normalization Idempotent]
\label{thm:normalization_idempotent}
\lean{InstitutionalClosure.normalization_idempotent}
\leanok
\uses{def:institutional_closure}
$\mathsf{closure} \circ \mathsf{temporalNormalize} = \mathsf{closure}$:
the normalization pipeline is a projection. Re-processing an already-
normalized history yields the same norm.
\end{theorem}

\section{Decision Composition API}

\begin{definition}[Gate]
\label{def:gate}
\lean{DecisionComposition.Gate}
\leanok
\uses{def:logiccontraction}
A Gate is a binary decision predicate indexed by a logic modality:
$\mathsf{Gate} = (\mathsf{name} : \mathsf{String}) \times
(\mathsf{modality} : \mathsf{LogicType}) \times
(\mathsf{check} : \mathsf{EMLTree} \to \mathsf{Prop})$.
Every logic type gives rise to a Gate via $\mathsf{ofLogicType}$,
and every tree passes any such gate (by \ref{thm:logic_contracts}).
\end{definition}

\begin{definition}[Decision]
\label{def:decision}
\lean{DecisionComposition.Decision}
\leanok
\uses{def:gate}
A Decision is the refinement type from the specification:
$\mathsf{Decision\ gates} = \{ d : \mathsf{EMLTree} \mid
\forall g \in \mathsf{gates},\ g.\mathsf{check}\ d \}$.
The type parameter $\mathsf{gates}$ IS the channel: downstream
functions know statically which gates the datum has passed.
\end{definition}

\begin{definition}[Compose]
\label{def:compose}
\lean{DecisionComposition.Decision.compose}
\leanok
\uses{def:decision}
$\mathsf{compose} : \mathsf{Decision\ gates} \to \mathsf{Gate}
\to \mathsf{Decision\ (g :: gates)}$.
The return type is strictly richer than the input type:
a gate is added to the type-level history.
\end{definition}

\begin{definition}[Left-Comb Normal Form]
\label{def:leftComb}
\lean{EMLRegistry.leftComb}
\leanok
\uses{def:emltree}
The mirror image of $\mathsf{rightComb}$: $\mathsf{leftComb}(0) = \mathsf{Leaf}$ and $\mathsf{leftComb}(n+1) = \mathsf{Node}(\mathsf{leftComb}(n),\ \mathsf{Leaf})$.
Unlike $\mathsf{rightComb}$, this is not the equilibrium attractor — it represents a fully left-associative configuration.
\end{definition}

\begin{definition}[LogicPipeline]
\label{def:logic_pipeline}
\lean{DecisionComposition.LogicPipeline}
\leanok
\uses{def:gate}
A LogicPipeline is a sequence of logic types to run as gates.
Running it on a tree produces a Decision tagged with all of them.
The institutional closure is a pipeline with gates Temporal,
Fuzzy, Deontic.
\end{definition}

\begin{theorem}[Pipeline Soundness]
\label{thm:pipeline_soundness}
\lean{DecisionComposition.closure_sound}
\leanok
\uses{def:logic_pipeline, def:decision}
For any tree and any gate in the closure pipeline,
$\mathsf{soundness}$ guarantees the gate check holds.
This is the compile-time contract: the type signature IS the channel.
\end{theorem}

\begin{theorem}[Decide Soundness]
\label{thm:decide_soundness}
\lean{DecisionComposition.decide_sound}
\leanok
\uses{def:logic_pipeline, def:decision, def:institutional_closure}
$\mathsf{decide}$ runs the institutional closure on $\mathsf{LogicM\ Event}$
and produces a Decision whose datum passes all closure gates.
The channel is the type: $\mathsf{Decision\ (closurePipeline.gates)}$.
\end{theorem}

\section{Paradox Taxonomy}

\begin{definition}[Problem Class]
\label{def:problemclass}
\lean{LiarParadox.ProblemClass}
\leanok
An inductive type enumerating 12 paradox classes, each with a native logic type:
self-reference (ManyValued), vagueness (Fuzzy), inconsistent definition (Paraconsistent),
temporal decision (Temporal), deontic (Deontic), epistemic (Epistemic),
quantum superposition (Quantum), constructive (Intuitionistic), relevance (Relevance),
empty reference (Free), infinity (Infinitary), modality (Modal).
\end{definition}

\begin{definition}[Problem]
\label{def:problem}
\lean{LiarParadox.Problem}
\leanok
\uses{def:problemclass, def:logictype, def:emltree}
A structure parameterized by a problem class, a name, a list of suitable logics,
a tree family $\mathsf{LogicType} \to \mathsf{EMLTree}$ (the paradox as encoded in each logic),
and a normal form family $\mathsf{LogicType} \to \mathsf{EMLTree}$ (the resolution target for each logic).
\end{definition}

\begin{definition}[Wrapped Problem]
\label{def:wrapped}
\lean{LiarParadox.WrappedProblem}
\leanok
\uses{def:problem, def:logiccontraction}
The WfCA collapse rule: pairs a Problem with a LogicType.
Records the tree (problem as interpreted in this logic), target (normal form),
cost (pentagonator distance $\Phi$), and proof (contraction path).
A problem's \textbf{superposition} of possible resolutions collapses to the
definite outcome given by the logic type's contraction rule.
\end{definition}

\section{The Liar as a Problem Instance}

\begin{definition}[Liar Problem]
\label{def:liar_problem}
\lean{LiarParadox.liarProblem}
\leanok
\uses{def:problem, def:problemclass}
The Liar Paradox stated as a \textsf{Problem} of class $\mathsf{selfReference}$.
Its tree is the size-3 symmetric tree $\mathsf{Node}(\mathsf{Node}(\mathsf{Leaf},\ \mathsf{Leaf}),\ \mathsf{Node}(\mathsf{Leaf},\ \mathsf{Leaf}))$
for most logics, except ManyValued which uses a size-2 left-comb encoding the
three-valued disjunction $\mathsf{True} \lor \mathsf{False} \lor \mathsf{Undefined}$.
All logics target $\mathsf{rightComb}(3)$ as the normal form.
\end{definition}

\section{Liar Wrappers}

\begin{definition}[Fuzzy Liar]
\label{def:fuzzy_liar}
\lean{LiarParadox.fuzzyLiar}
\leanok
\uses{def:wrapped, def:liar_problem}
The Liar wrapped in Fuzzy logic. The fixed point $\neg 0.5 = 0.5$ gives
pentagonator distance $\Phi = 1$ (CD step 1). Proven via the generic
\hyperref[def:liar_wrapper]{liarWrapper}.
\end{definition}

\begin{definition}[Meta-Paradox]
\label{def:metaparadox}
\lean{LiarParadox.missingProofParadox}
\leanok
\uses{def:problemclass, def:problem}
A \textsf{Problem} of class $\mathsf{metaParadox}$ constructed from any
existing Problem and LogicType whose proof is missing.
\end{definition}

\begin{definition}[Classical Liar]
\label{def:classical_liar}
\lean{LiarParadox.classicalLiar}
\leanok
\uses{def:wrapped, def:liar_problem, thm:contracts_to_rightComb}
The Liar wrapped in Classical logic. Cost $\Phi = 0$ (at equilibrium).
Proven using $\mathsf{contracts\_to\_rightComb}$: the symmetric tree contracts
to $\mathsf{rightComb}(3)$ in finitely many steps.
\end{definition}

\section{Nested Wrappers (Stacked Collapse)}

\begin{definition}[Tower]
\label{def:tower}
\lean{LiarParadox.Tower}
\leanok
\uses{def:wrapped}
A sequence of logic-wrapped problems, each collapsing the previous layer's output.
Implements the nested wrapper pattern from Lumo section B:
\[
\begin{array}{c}
\text{Layer 3 (Quantum): liar\_superposition} \\
\downarrow \\
\text{Layer 2 (Intuitionistic): proof\_of\_quantum\_state} \\
\downarrow \\
\text{Layer 1 (Fuzzy): confidence\_in\_proof} \\
\downarrow \\
\text{Layer 0 (Classical): assertion\_value}
\end{array}
\]
Each layer's logic type is tracked via a dependent pair $\Sigma lt,\ \mathsf{WrappedProblem}\ p\ lt$.
The total cost is the sum of per-layer costs.
\end{definition}

\begin{definition}[Liar Tower]
\label{def:liar_tower}
\lean{LiarParadox.liarTower}
\leanok
\uses{def:tower, def:classical_liar, def:fuzzy_liar}
A four-layer prototype tower alternating Classical and Fuzzy wrappers.
Demonstrates the stacked collapse pattern; the full tower with Quantum,
Intuitionistic, Fuzzy, and Classical layers is not yet populated because
non-Classical \textsf{LogicContraction}s are missing.
\end{definition}

\section{Temporal Paradox Problem}

\begin{definition}[Grandfather Problem]
\label{def:grandfather_problem}
\lean{TemporalParadox.grandfatherProblem}
\leanok
\uses{def:problem, def:problemclass}
The Grandfather Paradox stated as a \textsf{Problem} of class $\mathsf{temporalDecision}$.
Its tree is a left-comb chain of size 4, encoding the 4-step time-travel causal loop:
traveler exists $\to$ travels back $\to$ grandfather killed $\to$ traveler never born (contradiction).
The normal form is $\mathsf{rightComb}(4)$.
\end{definition}

\begin{definition}[Generic Grandfather Wrapper]
\label{def:grandfather_wrapper}
\lean{TemporalParadox.grandfatherWrapper}
\leanok
\uses{def:grandfather_problem, def:wrapped}
For any logic \textsf{lt}, constructs the \textsf{WrappedProblem} with
$\textsf{cost} = \textsf{cdStep}(\textsf{lt})$. Resolves the temporal causal loop
via $\mathsf{contracts\_to\_rightComb}$.
\end{definition}

\begin{definition}[Grandfather Tower]
\label{def:grandfather_tower}
\lean{TemporalParadox.grandfatherTower}
\leanok
\uses{def:grandfather_wrapper, def:tower}
The Grandfather tower: one layer per suitable logic.
\end{definition}

\begin{definition}[Grandfather Cost]
\label{def:grandfather_cost}
\lean{TemporalParadox.grandfatherCost}
\leanok
\uses{def:cdstep}
The cost $\Phi(\text{Grandfather}, \text{lt})$, equal to $\textsf{cdStep}(\text{lt})$.
\end{definition}

\begin{theorem}[Grandfather Cost Bound]
\label{thm:grandfather_cost_bound}
\lean{TemporalParadox.grandfatherCost_le_cdStep}
\leanok
\uses{def:grandfather_cost}
$\textsf{grandfatherCost}(\text{lt}) \le \textsf{cdStep}(\text{lt})$ for every logic type.
\end{theorem}

\section{Russell's Paradox Problem}

\begin{definition}[Russell's Problem]
\label{def:russells_problem}
\lean{RussellsParadox.russellsProblem}
\leanok
\uses{def:problem, def:problemclass}
Russell's Paradox stated as a \textsf{Problem} of class $\mathsf{inconsistentDef}$.
Its tree is a left-comb chain of size 3, encoding the three levels of
set-theoretic diagonalization: elements $\to$ self-membership predicate $\to$
the set of all sets not containing themselves.
The normal form is $\mathsf{rightComb}(3)$.

\textbf{Disclaimer:} This formalization is deliberately incomplete with respect to
the infinity/eternity axis. The ``cheat'' of setting cost $= \textsf{cdStep}$
captures only the Cayley--Dickson property-loss dimension, not the
regularization dimension (endless $\to$ eternity $\to$ actual $\infty \to$
boundlessness). Future work must pin this axis down properly; here we
procrastinate it using the same wrapper pattern as Liar/Sorites/Grandfather.
\end{definition}

\begin{definition}[Generic Russell Wrapper]
\label{def:russells_wrapper}
\lean{RussellsParadox.russellsWrapper}
\leanok
\uses{def:russells_problem, def:wrapped}
For any logic \textsf{lt}, constructs the \textsf{WrappedProblem} with
$\textsf{cost} = \textsf{cdStep}(\textsf{lt})$. Resolves set-theoretic
diagonalization via $\mathsf{contracts\_to\_rightComb}$.
\end{definition}

\begin{definition}[Russell Tower]
\label{def:russells_tower}
\lean{RussellsParadox.russellsTower}
\leanok
\uses{def:russells_wrapper, def:tower}
The Russell tower: one layer per suitable logic.
\end{definition}

\begin{definition}[Russell Cost]
\label{def:russells_cost}
\lean{RussellsParadox.russellsCost}
\leanok
\uses{def:cdstep}
The cost $\Phi(\text{Russell}, \text{lt})$, equal to $\textsf{cdStep}(\text{lt})$.
\end{definition}

\begin{theorem}[Russell Cost Bound]
\label{thm:russells_cost_bound}
\lean{RussellsParadox.russellsCost_le_cdStep}
\leanok
\uses{def:russells_cost}
$\textsf{russellsCost}(\text{lt}) \le \textsf{cdStep}(\text{lt})$ for every logic type.
\end{theorem}

\section{The Very Big Box}

The Very Big Box is the product space
\[
\mathcal{V} = \prod_{c \in \mathsf{ProblemClass}}\ \prod_{lt \in \mathsf{suitableLogics}(c)}
\mathsf{WrappedProblem}(\mathsf{problem}(c),\ lt)
\]
where each coordinate is a problem in a specific logic and the
Tower construction provides the nesting structure. Currently four
problem classes are populated: $\mathsf{selfReference}$ (Liar),
$\mathsf{vagueness}$ (Sorites), $\mathsf{temporalDecision}$ (Grandfather),
and $\mathsf{inconsistentDef}$ (Russell's).
The remaining 9 classes represent the \textbf{real friction}: every cell
in this product that lacks a proof is a concrete blocker.

\section{Sorites as a Problem Instance}

\begin{definition}[Sorites Problem]
\label{def:sorites_problem}
\lean{SoritesParadox.soritesProblem}
\leanok
\uses{def:problem, def:problemclass}
The Sorites (Heap) Paradox stated as a \textsf{Problem} of class $\mathsf{vagueness}$.
Its tree is a left-comb chain of size 5, encoding the inductive threshold:
``If $n$ grains form a heap, then $n-1$ grains form a heap.''
The normal form is $\mathsf{rightComb}(5)$.
\end{definition}

\begin{definition}[Generic Sorites Wrapper]
\label{def:sorites_wrapper}
\lean{SoritesParadox.soritesWrapper}
\leanok
\uses{def:sorites_problem, def:wrapped, def:liar_wrapper}
For any logic \textsf{lt}, constructs the \textsf{WrappedProblem} with
$\textsf{cost} = \textsf{cdStep}(\textsf{lt})$. Follows the same pattern
as the generic Liar wrapper, using $\mathsf{contracts\_to\_rightComb}$.
\end{definition}

\begin{definition}[Sorites Tower]
\label{def:sorites_tower}
\lean{SoritesParadox.soritesTower}
\leanok
\uses{def:sorites_wrapper, def:tower}
The Sorites tower: one layer per suitable logic, analogous to the Liar tower.
\end{definition}

\begin{definition}[Sorites Cost]
\label{def:sorites_cost}
\lean{SoritesParadox.soritesCost}
\leanok
\uses{def:cdstep}
The cost $\Phi(\text{Sorites}, \text{lt})$, equal to $\textsf{cdStep}(\text{lt})$.
\end{definition}

\begin{theorem}[Sorites Cost Bound]
\label{thm:sorites_cost_bound}
\lean{SoritesParadox.soritesCost_le_cdStep}
\leanok
\uses{def:sorites_cost}
$\textsf{soritesCost}(\text{lt}) \le \textsf{cdStep}(\text{lt})$ for every logic type.
\end{theorem}

\section{Decomposition: Ontological Underdetermination}

\begin{definition}[Path]
\label{def:path}
\lean{Decomposition.Path}
\leanok
\uses{def:contracts_one}
A witness of a contraction path through the Tamari lattice:
$\mathsf{Path}\ s\ t$ is an inductive sequence of
$\mathsf{contracts\_one}$ steps from $s$ to $t$.
\begin{itemize}
  \item $\mathsf{nil}$ --- the trivial path (reflexivity).
  \item $\mathsf{cons}\ h\_one\ p$ --- prepend a single contraction step to an existing path.
\end{itemize}
\end{definition}

\begin{theorem}[Path Soundness]
\label{thm:path_soundness}
\lean{Decomposition.Path.to_contracts_to}
\leanok
\uses{def:path, def:contracts_to}
For any path $p : \mathsf{Path}\ s\ t$, we have $\mathsf{contracts\_to}\ s\ t$.
The path provides constructive evidence for the existence of a contraction sequence.
\end{theorem}

\begin{definition}[Path Length]
\label{def:path_length}
\lean{Decomposition.Path.length}
\leanok
\uses{def:path}
The number of contraction steps in a path:
$\mathsf{length}(\mathsf{nil}) = 0$,
$\mathsf{length}(\mathsf{cons}\ h\_one\ p) = 1 + \mathsf{length}(p)$.
\end{definition}

\begin{definition}[Path Concatenation]
\label{def:path_append}
\lean{Decomposition.Path.append}
\leanok
\uses{def:path}
Concatenates two compatible paths: $\mathsf{append} : \mathsf{Path}\ s\ t \to \mathsf{Path}\ t\ u \to \mathsf{Path}\ s\ u$.
\end{definition}

\section{Reverse Contraction: Generating Predecessors}

\begin{definition}[Reverse One]
\label{def:reverse_one}
\lean{Decomposition.reverse_one}
\leanok
\uses{def:contracts_one, def:emltree}
Generates all immediate predecessors of a tree under $\mathsf{contracts\_one}$.
For a target tree $t$, $\mathsf{reverse\_one}(t)$ is a list of all $s$ such that
$s \to t$ in one step. Three structural cases:
\begin{itemize}
  \item $\mathsf{Leaf}$ has no predecessors.
  \item $\mathsf{Node}\ l\ (\mathsf{Node}\ r_1\ r_2)$: the rotation target plus
    predecessors from reversing the left and right subtrees.
  \item $\mathsf{Node}\ l\ r$ (right child not a Node): predecessors from
    reversing the left and right subtrees only.
\end{itemize}
This is the core de-composition primitive: given an outcome, what could have
immediately preceded it?
\end{definition}

\begin{theorem}[Reverse One Soundness]
\label{thm:reverse_one_sound}
\lean{Decomposition.reverse_one_sound}
\leanok
\uses{def:reverse_one, def:contracts_one}
If $s \in \mathsf{reverse\_one}(t)$ then $\mathsf{contracts\_one}\ s\ t$.
Every generated predecessor is a valid immediate prior configuration.
\end{theorem}

\begin{theorem}[Reverse One Completeness]
\label{thm:reverse_one_complete}
\lean{Decomposition.reverse_one_complete}
\leanok
\uses{def:reverse_one, def:contracts_one}
If $\mathsf{contracts\_one}\ s\ t$ then $s \in \mathsf{reverse\_one}(t)$.
Every valid immediate predecessor is generated.
\end{theorem}

\section{Decomposition Structure}

\begin{definition}[Decomposition]
\label{def:decomposition}
\lean{Decomposition.Decomposition}
\leanok
\uses{def:contracts_to}
For a target tree $\mathsf{target}$, a $\mathsf{Decomposition}$ pairs a
$\mathsf{source}$ with a proof $\mathsf{contracts\_to}\ \mathsf{source}\ \mathsf{target}$.
This is the \textbf{ontological unit} of de-composition: an outcome together with
a specific prior configuration that could have produced it.
\end{definition}

\begin{theorem}[Left-Comb Size]
\label{thm:leftcomb_size}
\lean{Decomposition.leftComb_size}
\leanok
\uses{def:rightComb}
$\mathsf{size}(\mathsf{leftComb}\ n) = n$.
\end{theorem}

\begin{theorem}[Right-Comb Size]
\label{thm:rightcomb_size}
\lean{Decomposition.rightComb_size}
\leanok
\uses{def:rightComb}
$\mathsf{size}(\mathsf{rightComb}\ n) = n$.
\end{theorem}

\begin{theorem}[Right-Comb vs Left-Comb]
\label{thm:rightcomb_ne_leftcomb}
\lean{Decomposition.rightComb_ne_leftComb}
\leanok
\uses{def:rightComb, def:leftComb}
For $n \ge 2$, $\mathsf{rightComb}\ n \neq \mathsf{leftComb}\ n$.
The two canonical trees of size $n$ are distinct in the Tamari lattice.
\end{theorem}

\begin{theorem}[Non-Unique Decomposition]
\label{thm:non_unique_decomposition}
\lean{Decomposition.non_unique_decomposition}
\leanok
\uses{def:decomposition, def:rightComb, thm:rightcomb_ne_leftcomb}
For $n \ge 2$, the target $\mathsf{rightComb}\ n$ has at least two distinct
Decompositions with different sources:
\[
\exists d_1\ d_2 : \mathsf{Decomposition}(\mathsf{rightComb}\ n),\ d_1.\mathsf{source} \neq d_2.\mathsf{source}.
\]
This is the core ontological result: \textbf{the past is underdetermined by the
present}. The same outcome (rightComb $n$) admits multiple evidential histories
(rightComb $n$ itself and leftComb $n$, among others).
\end{theorem}

\begin{theorem}[Path Diversity]
\label{thm:path_diversity}
\lean{Decomposition.path_diversity}
\leanok
\uses{def:path, def:contracts_one, thm:rightcomb_ne_leftcomb}
There exist two distinct paths between the same source-target pair in $T_3$:
\[
\exists s\ t\ (p_1\ p_2 : \mathsf{Path}\ s\ t),\ p_1 \neq p_2.
\]
Explicitly, $\mathsf{leftComb}\ 3 \to \mathsf{rightComb}\ 3$ has two
distinct maximal chains in the Tamari lattice:
\begin{itemize}
  \item $((ab)c)d \to (a(bc))d \to a((bc)d) \to a(b(cd))$ (3 steps)
  \item $((ab)c)d \to (ab)(cd) \to a(b(cd))$ (2 steps)
\end{itemize}
This proves that \textbf{intent is not uniquely determined by outcome} ---
multiple reasoning chains can connect the same evidence to the same verdict.
\end{theorem}

\section{The Hypercomputer: Infinite Ancestor Enumeration}

\begin{definition}[Predecessors]
\label{def:predecessors}
\lean{Decomposition.predecessors}
\leanok
\uses{def:reverse_one}
The immediate predecessor sources of a target tree $\mathsf{predecessors}(t) = \mathsf{reverse\_one}(t)$.
Soundness is maintained as a separate theorem $\mathsf{predecessors\_sound}$:
if $s \in \mathsf{predecessors}(t)$ then $\mathsf{contracts\_one}\ s\ t$.
\end{definition}

\begin{definition}[Chain]
\label{def:chain}
\lean{Decomposition.Chain}
\leanok
\uses{def:contracts_one}
A finite chain of contraction steps from an ancestor to a target.
$\mathsf{Chain.tip}\ t$ is the empty chain (the present moment without a
reconstructed past). $\mathsf{Chain.link}\ h\ r$ prepends a contraction step.
\end{definition}

\begin{definition}[Ancestors Up To Depth $n$]
\label{def:ancestorsUpTo}
\lean{Decomposition.ancestorsUpTo}
\leanok
\uses{def:predecessors}
Enumerate all ancestor sources up to depth $n$ via DFS, returning a flat list.
This is the \textbf{hypercomputer in function form}: the infinite tree of all
possible pasts is present as the limit of $\mathsf{ancestorsUpTo}\ t\ n$ as
$n \to \infty$. No single data structure holds the whole tree; instead the
recursion is explicit in the partial function, and every finite approximation
is immediately available. ``Instant lookup'' = calling this function with
any depth $n$.
\end{definition}

\begin{definition}[DFS View]
\label{def:viewDFS}
\lean{Decomposition.viewDFS}
\leanok
\uses{def:chain, def:predecessors}
A $\mathsf{viewDFS}$ is a single lineage (DFS) extracted from the ancestor
space by always following the first predecessor. Returns a $\mathsf{Chain}$
that may end at any depth. This is the \textbf{narrative self} in formal
type: committed to a single thread, but honest about termination.
\end{definition}

\paragraph{Lean 4 Coinductive Limitation}
\lean{Decomposition.lean4_limitation_note}
\leanok
The infinite ancestor tree cannot be represented as a closed coinductive
data type in Lean 4 because \texttt{coinductive} is restricted to
\texttt{Prop}-valued predicates, and nested inductives with \texttt{Prop}
fields through \texttt{List} and \texttt{Sigma} are rejected by the kernel
positivity checker. The workaround (partial functions) is not a limitation
of the framework but a \textbf{design choice}: the hypercomputer exists as
the limit of finite approximations, not as a closed value. This matches
the ontology: the infinite crystal seed is not a thing we hold, but a
thing we approach.

\section{Boundlessness: The Terminal Regularization}
\label{sec:boundlessness}

The regularization ladder (endless $\to$ eternal $\to$ actual $\infty \to$
boundlessness) is capped by the \textbf{idempotent resolution} framework:
a boundary-response operation that is already idempotent is already at its
own limit.

\begin{definition}[Idempotent Resolution]
\label{def:idempotent_resolution}
\lean{Boundlessness.IdempotentResolution}
\lean{Boundlessness.rightCombResolution}
\leanok
\uses{def:rightComb, def:size}
An \textbf{idempotent resolution} for a type $\alpha$ consists of:
\begin{itemize}
  \item $\mathsf{step} : \alpha \to \alpha$, a regularization step that is
    idempotent ($\mathsf{step} \circ \mathsf{step} = \mathsf{step}$).
  \item $\mathsf{limit} : \alpha \to \alpha$, the limiting configuration
    reachable in one step from any image of $\mathsf{step}$
    ($\mathsf{limit} = \mathsf{step} \circ \mathsf{limit}$).
\end{itemize}
The canonical example is \textsf{rightCombResolution}, where
$\mathsf{step}(t) = \mathsf{rightComb}(t.\mathsf{size})$: every tree is
mapped to its right-comb normal form in one step, and applying the step
again does nothing because a right-comb is already in normal form.
\end{definition}

\begin{theorem}[Right-Comb Size]
\label{thm:rightcomb_size_boundlessness}
\lean{Boundlessness.rightComb_size}
\leanok
\uses{def:rightComb}
$\mathsf{size}(\mathsf{rightComb}\ n) = n$ for all $n : \mathbb{N}$.
\end{theorem}

\begin{definition}[Very Big Box]
\label{def:very_big_box}
\lean{Boundlessness.VeryBigBox}
\lean{Boundlessness.veryBigBox}
\leanok
\uses{def:idempotent_resolution}
The \textbf{Very Big Box} is the product of idempotent resolutions over all
populated problem classes:
\[
\mathcal{V} = \prod_{c \in \{\text{Liar, Sorites, Grandfather, Russell}\}}
\mathsf{IdempotentResolution}\ \mathsf{EMLTree}
\]
The canonical \textsf{veryBigBox} assigns \textsf{rightCombResolution}
to each coordinate.
\end{definition}

\begin{theorem}[Meta-Idempotence]
\label{thm:rightComb_meta_idemp}
\lean{Boundlessness.rightComb_meta_idemp}
\leanok
\uses{def:idempotent_resolution}
The step function of \textsf{rightCombResolution} is already its own fixed
point:
\[
\mathsf{rightCombResolution.step} \circ \mathsf{rightCombResolution.step}
= \mathsf{rightCombResolution.step}
\]
This is the terminal regularization: applying boundary-response again
changes nothing. Not because there is nothing more to say, but because
saying more is already captured by the first application.
\end{theorem}

\begin{theorem}[Limit Idempotence]
\label{thm:rightComb_limit_idemp}
\lean{Boundlessness.rightComb_limit_idemp}
\leanok
\uses{def:rightComb, def:idempotent_resolution}
Applying \textsf{rightComb} normal form to a tree that is already a
\textsf{rightComb} is the identity:
\[
\mathsf{rightComb}((\mathsf{rightComb}\ t.\mathsf{size}).\mathsf{size}) = \mathsf{rightComb}\ t.\mathsf{size}
\]
Boundlessness is not a structure we build, but a property we recognize.
\end{theorem}

\paragraph{Relation to the Hypercomputer.}
The idempotent resolution caps the hypercomputer: where $\mathsf{ancestorsUpTo}$
approaches the infinite ancestor tree by finite approximation,
the idempotent resolution declares that the limit has been reached.
The Very Big Box is the product closure of all problem classes, and
boundlessness is the fixed point that requires no further iteration.
This is the formal declaration that the framework is closed under its
own boundary-response operation.
\section{Completed Work}
\label{sec:completed}

All \textsf{LogicContraction}, \textsf{LogicNormalForm}, and
\textsf{logic\_contracts\_to\_normal\_form} definitions are filled and proven
for all 13 logic types. The Liar problem has a complete \textsf{WrappedProblem}
for each suitable logic via \textsf{liarWrapper}. The Sorites and Grandfather
problems follow the same pattern with \textsf{soritesWrapper} and
\textsf{grandfatherWrapper}.

The Decomposition infrastructure (\textsf{Path}, \textsf{reverse\_one},
\textsf{Decomposition}, \texttt{predecessors}, \texttt{ancestorsUpTo},
\texttt{Chain}, \texttt{viewDFS}) is complete with soundness, completeness,
non-unique-decomposition, and path-diversity theorems proven for all
$n \ge 2$ in $T_n$ — zero sorries.

The Friction Lagrangian $\mathcal{L}_{\text{friction}}$ is computed from the
full Liar tower. The blocking chain for the Liar is fully resolved.

\subsection{Remaining Work}

\begin{itemize}
  \item \textbf{Other problem classes} --- $\mathsf{selfReference}$ (Liar),
    $\mathsf{vagueness}$ (Sorites), $\mathsf{temporalDecision}$ (Grandfather),
    and $\mathsf{inconsistentDef}$ (Russell's) are instantiated.
    The other 9 classes (deontic, epistemic, quantum, etc.)
    need \textsf{Problem} definitions with their own tree families.
  \item \textbf{Tower composition} --- the current towers are flat (all logics stacked).
    The institutional closure pipeline demonstrates a first step toward
    nested composition (temporal $\to$ fuzzy $\to$ deontic $\to$ self-recognition).
    Full \textsf{LogicTranslation} instances between adjacent logics remain.
  \item \textbf{Shortest-path cost} --- the $\Phi$ cost currently equals
    $\textsf{cdStep}$ by definition; a BFS algorithm over $\textsf{reverse\_one}$
    would compute the actual PentagonatorDistance.
  \item \textbf{Ancestor enumeration} --- $\mathsf{reverse\_one}$ can be iterated
    to produce $\mathsf{ancestorsBounded}$, a bounded BFS generating all distinct
    sources for a given target with $\mathsf{Path}$ witnesses.
  \item \textbf{Path cost} --- a $\mathsf{pathCost}\ lt\ p$ function summing
    $\mathsf{cdStep}$ per logic type along a path, enabling cross-logic cost
    comparison for the same reasoning chain.
  \item \textbf{Decision decomposition} --- extend the DecisionComposition API
    with $\mathsf{decisionDecompositions}$ that enumerates all prior decisions
    that could compose to a given decision.
  \item \textbf{Regularization axis} --- the institutional closure sketch
    uses placeholder normalization; the full endless/eternal/infinite/boundless
    hierarchy must be formalized to compute real costs.
\end{itemize}

\subsection{Liar Cost: Hardness Measure}

The Liar as perfect anti-coherence ($X = \neg X$) gives a hardness measure for each logic:
how many contraction steps are required to resolve it.

\begin{definition}[Liar Cost]
\label{def:liar_cost}
\lean{LiarParadox.liarCost}
\leanok
\uses{def:liar_problem}
The cost $\Phi(\text{Liar}, \text{lt})$ to resolve the Liar under logic \textsf{lt}.
Conjectured to equal the Cayley--Dickson step.
\end{definition}

\begin{theorem}[Cost Bound]
\label{thm:liar_cost_bound}
\lean{LiarParadox.liarCost_le_cdStep}
\leanok
\uses{def:liar_cost, def:cdstep}
$\textsf{liarCost}(\text{lt}) \le \textsf{cdStep}(\text{lt})$ for every logic type.
For Classical and Fuzzy the bound is tight (equality).
\end{theorem}

\begin{definition}[Generic Liar Wrapper]
\label{def:liar_wrapper}
\lean{LiarParadox.liarWrapper}
\leanok
\uses{def:liar_problem, def:liar_cost}
For any logic \textsf{lt}, constructs the \textsf{WrappedProblem} with
$\textsf{cost} = \textsf{liarCost}(\text{lt})$. Fills all 12 mail slots.
\end{definition}

\begin{corollary}[Cost Match for Classical]
\label{cor:liar_cost_match_classical}
\lean{LiarParadox.liarCost_matches_classical}
\leanok
\uses{thm:liar_cost_bound}
$\textsf{liarCost}(\text{Classical}) = 0 = \textsf{classicalLiar.cost}$.
\end{corollary}

\begin{corollary}[Cost Match for Fuzzy]
\label{cor:liar_cost_match_fuzzy}
\lean{LiarParadox.liarCost_matches_fuzzy}
\leanok
\uses{thm:liar_cost_bound}
$\textsf{liarCost}(\text{Fuzzy}) = 1 = \textsf{fuzzyLiar.cost}$.
\end{corollary}

\subsection{Friction Lagrangian}

The tower sums costs across all logic layers. The total \textbf{Friction Lagrangian}
$\mathcal{L}_{\text{friction}} = \sum_{\text{lt}} \Phi(\text{Liar}, \text{lt})$ is
the total resistance of the logical ecosystem to self-negation.

\begin{definition}[Friction Lagrangian]
\label{def:friction_lagrangian}
\lean{LiarParadox.frictionLagrangian}
\leanok
\uses{def:liar_tower, def:liar_cost}
The sum of $\Phi$ costs across all layers of the Liar tower.
\end{definition}

With all 12 suitable logics, the Friction Lagrangian converges to a stable value:
\[
\mathcal{L}_{\text{friction}} = \sum_{\text{lt} \in \text{suitableLogics}} \textsf{cdStep}(\text{lt})
\]

This is the ``curvature of mind'' --- the total structural work the system must do
to resolve perfect anti-coherence across the full space of logics.

The blocking chain is now fully resolved for the Liar problem.
Previously the sorries in \textsf{LogicTypes.lean} were abstract goals;
now every definition and theorem is filled, and the Very Big Box has its
first complete product cell: the Liar across all 12 logics.

\endinput